\section{Review Exercises}

\subsection{Elementary computational and construction exercises}

\begin{enumerate}
\item Write a parametric equation of the line through $P=(2,-1)$ with direction vector $v=(3,4)$. Find the points corresponding to $t=-1,0,2$ and plot them.

\item Find an implicit equation of the line through $P=(1,3)$ with normal vector $n=(2,-1)$. Check directly that $P$ satisfies the equation.

\item Convert the line
\[
2x-3y=6
\]
to slope-intercept form, find both axis intercepts, and give one direction vector and one normal vector.

\item Find the equation of the line through
\[
A=(-1,2),\qquad B=(3,-4)
\]
in parametric form and, when possible, in slope form.

\item Find the intersection point of
\[
L_1:\ y=2x-1,
\qquad
L_2:\ x+y=5.
\]
Verify the point in both equations.

\item Determine whether each pair of lines is parallel, perpendicular, identical, or intersecting without either special relation:
\begin{enumerate}
\item $2x-y=1$ and $4x-2y=7$;
\item $x+2y=3$ and $2x-y=4$;
\item $3x+y=5$ and $6x+2y=10$.
\end{enumerate}

\item Compute the distance from $P=(3,1)$ to the line
\[
2x-y-4=0.
\]
Find the foot of the perpendicular projection and verify the distance using coordinates.

\item In space, write a parametric equation of the line through
\[
P=(1,0,-2)
\]
with direction $v=(2,-1,3)$. Test whether $Q=(5,-2,4)$ lies on the line.

\item Write a parametric equation of the plane through
\[
P=(1,2,0)
\]
with direction vectors $u=(1,0,1)$ and $v=(0,2,-1)$. Find three distinct points of the plane.

\item Find a normal vector and an implicit equation of the plane through $P=(1,-1,2)$ spanned by
\[
u=(1,1,0),\qquad v=(0,1,1).
\]
Check that both direction vectors are perpendicular to the normal.
\end{enumerate}

\subsection{Conceptual and reasoning exercises}

\begin{enumerate}
\item Explain why one point and one nonzero direction determine a line, while one point alone does not.

\item Compare parametric, normal, implicit, and slope descriptions of a line. For each form, state one kind of question that it answers especially efficiently.

\item Why does slope notation require a separate vertical-line case while the parametric description does not? Explain what this reveals about representations and geometric objects.

\item Prove that the line through distinct points $P$ and $Q$ can be written as
\[
L=\{P+t(Q-P):t\in\mathbb R\}.
\]
Explain why every point in this set lies on the same geometric line.

\item Two equations look different but may describe the same line. Give a criterion for recognising this in implicit form and justify it.

\item Explain why finding an intersection means imposing two defining conditions simultaneously. Relate this principle to axis intercepts and to the intersection of two arbitrary lines.

\item Show why equal nonzero normal vectors up to scaling produce parallel or identical lines. What additional information distinguishes the two cases?

\item Explain why one point and two nonparallel directions determine a plane. What geometric object is obtained if the two directions are parallel?

\item Compare a plane described by two spanning directions with a plane described by a normal vector. Why are these two kinds of information complementary?

\item A student solves every line problem by converting immediately to slope form. Explain why this strategy can hide structure or create unnecessary exceptional cases. Give two examples where another representation is preferable.
\end{enumerate}
